\scanpage{203}
\mspage{749}

\uline{ReN.B.} Unless I am mistaken, property \(1^\circ)\) means that one can find an affine group scheme \(G\) of finite type over \(\mathbb{Z}\), which \unsure{defines|defines} \(G_{\mathbb{Q}}\), and such that (27) factors through \(G(\widehat{\mathbb{Z}})\).
\diagram{\begin{tikzcd}[ampersand replacement=\&]
(29)\quad \deleted{\unsure{G}}\quad \mathcal{M} \arrow[r] \& G(\widehat{\mathbb{Z}}).
\end{tikzcd}}

\(2^\circ)\) \uline{Compatibility with} \(\chi\) :

\(\unsure{(30)}\)
\diagram{\begin{tikzcd}[ampersand replacement=\&, row sep=large, column sep=large]
\& \mathcal{M} \arrow[r] \arrow[dl, "\chi_{\mathcal{M}}"'] \& G(\mathcal{A}_{\mathbb{Q}}) \arrow[dl, "\chi_G"] \\
\widehat{\mathbb{Z}} \arrow[r, hook] \& \unsure{\mathbb{G}_m|G_m}(\mathcal{A}_{\mathbb{Q}})\simeq\mathcal{A}_{\mathbb{Q}}^{*}\subset\displaystyle\prod_{\ell\in\mathbb{P}}\mathbb{Q}_{\ell}^{*} \&
\end{tikzcd}}
is commutative.

\(3^\circ)\). \uline{Compatibility with the discretification} \(\Pi_0\).

\((31)\)
\diagram{\begin{tikzcd}[ampersand replacement=\&, row sep=large, column sep=large]
\Pi'_0 \arrow[r, hook] \arrow[d, "\text{factorization}"] \& \Pi \arrow[r, hook] \& \mathcal{M} \arrow[d] \\
G(\mathbb{Q}) \arrow[rr, hook] \&\& G(\mathcal{A}_{\mathbb{Q}})
\end{tikzcd}},
for a suitable subgroup of finite index \(\Pi'_0\) of \(\Pi_0\).

\uline{N.B.} Changing notation if need be \unsure{we may assume|we may assume} \(\Pi'_0=\Pi_0\).

\(4^\circ)\) \uline{Minimality} (\unsure{irredundancy|irredundancy}) :

If \(G'_{\mathbb{Q}}\) is a subgroup scheme of

\scanpage{204}

\(G_{\mathbb{Q}}\), such that \(\deleted{\unsure{G}}\ \mathcal{M}\longrightarrow G(\mathcal{A}_{\mathbb{Q}})\) factors through \(G'(\mathcal{A}_{\mathbb{Q}})\), then \(G=G'\).

\uline{N.B.} If this hypothesis is not satisfied, then among all the \(G'_{\mathbb{Q}}\subset G_{\mathbb{Q}}\) which violate it, there is a \inserted{(unique)} smallest one, and on replacing \(\chi_G\) by \(\chi_{G'}=\chi_G\mid G'_{\mathbb{Q}}\ldots\), the hypotheses \(1^\circ)\) to \(4^\circ)\) are satisfied for \(G'_{\mathbb{Q}},\chi'\), and \(\mathcal{M}\longrightarrow G'(\mathcal{A}_{\mathbb{Q}})\).

\uline{Remark}. Let \(\mathcal{G}_0\subset\mathcal{M}\) be the subgroup of \(\mathcal{M}\) formed of the \(g\in\mathcal{M}\) which are such that \(\unsure{g(\Pi_0)|g\ell(\Pi_0)}\) and \(\Pi_0\) are commensurable. \unsure{That is|That is} \inserted{the \(g\)} which normalize the ``discretification germ'' \(\Pi_0\) of \(\Pi\).

\margin{\(\mathcal{G}_0/\unsure{\Pi_0}\) plays the role of a \uline{discrete} Teichmüller group}

It is possible that condition \(3^\circ\) will have to be strengthened, by requiring that we have
\diagram{\begin{tikzcd}[ampersand replacement=\&]
\mathcal{G}_0 \arrow[r] \& G(\mathbb{Q}).
\end{tikzcd}}
The future will decide it!

\scanpage{205}
\mspage{750}

Reflection \unsure{done|done}, I prefer to put things in a slightly \unsure{different|different} way, by introducing, in the \inserted{profinite} \(LT\)-group \(\mathcal{M}\), a normal \inserted{subgroup} \deleted{defining the \(LT\)-\unsure{structure|structure}} \(\Pi\) with an \(L\)-structure (not \unsure{necessari\textsuperscript{ly}|necessarily} a loop structure --- for example one can have \(\Pi=\mathrm{SL}\unsure{(}2,\mathbb{Z})\)) defining the \(LT\)-structure of \(\mathcal{M}\), a discretification \(\Pi_0\) of \(\Pi\) (in an obvious sense, for \(L\)-structures) a discrete subgroup \(\mathcal{G}_0\) \inserted{\(\supset\Pi_0\)} normalizing \(\Pi_0\) (N.B. in \unsure{important} special cases \unsure{will be} \(\mathcal{G}_0=\Pi_0\)), such that \inserted{its closure} \(\mathcal{G}_0^{\deleted{\unsure{\wedge|*}}}\) \inserted{\(\overset{\mathrm{dfn}}{=}\mathcal{G}\)} in \(\mathcal{M}\) is normal in \(\mathcal{M}\).

We thus have a diagram of groups

\((32)\)
\diagram{\begin{tikzcd}[ampersand replacement=\&,
row sep=large, column sep=large,
execute at end picture={
  \path (s205p.north) -- coordinate[midway] (s205pgmid) (s205g.north);
  \path (s205g.north) -- coordinate[midway] (s205gmmid) (s205m.north);
  \path (s205p.north) -- coordinate[midway] (s205pmmid) (s205m.north);
  \path (s205pzero.south) -- coordinate[midway] (s205zeromid) (s205gzero.south);
  \draw ([yshift=5mm]s205p.north) .. controls ([yshift=8mm]s205p.north) and ([xshift=-2mm,yshift=5mm]s205pgmid) .. ([yshift=8mm]s205pgmid) .. controls ([xshift=2mm,yshift=5mm]s205pgmid) and ([yshift=8mm]s205g.north) .. ([yshift=5mm]s205g.north);
  \node[above] at ([yshift=8mm]s205pgmid) {$\mathcal{T}$};
  \draw ([yshift=5mm]s205g.north) .. controls ([yshift=8mm]s205g.north) and ([xshift=-2mm,yshift=5mm]s205gmmid) .. ([yshift=8mm]s205gmmid) .. controls ([xshift=2mm,yshift=5mm]s205gmmid) and ([yshift=8mm]s205m.north) .. ([yshift=5mm]s205m.north);
  \node[above] at ([yshift=8mm]s205gmmid) {$\Gamma$};
  \draw ([yshift=17mm]s205p.north) .. controls ([yshift=20mm]s205p.north) and ([xshift=-2mm,yshift=17mm]s205pmmid) .. ([yshift=20mm]s205pmmid) .. controls ([xshift=2mm,yshift=17mm]s205pmmid) and ([yshift=20mm]s205m.north) .. ([yshift=17mm]s205m.north);
  \node[above] at ([yshift=20mm]s205pmmid) {$\mathcal{N}$};
  \draw ([yshift=-3mm]s205pzero.south) .. controls ([yshift=-6mm]s205pzero.south) and ([xshift=-2mm,yshift=-3mm]s205zeromid) .. ([yshift=-6mm]s205zeromid) .. controls ([xshift=2mm,yshift=-3mm]s205zeromid) and ([yshift=-6mm]s205gzero.south) .. ([yshift=-3mm]s205gzero.south);
  \node[below] at ([yshift=-6mm]s205zeromid) {$\mathcal{T}_0$};
}]
\unsure{1} \arrow[r, hook] \& |[alias=s205p]| \underset{=\widehat{\Pi}_0}{\Pi} \arrow[r, hook] \& |[alias=s205g]| \underset{=\widehat{\mathcal{G}}_0}{\mathcal{G}^{\deleted{\unsure{\wedge|*}}}_{\deleted{\unsure{0}}}} \arrow[r, hook] \& |[alias=s205m]| \mathcal{M} \\
1 \arrow[r, hook] \& |[alias=s205pzero]| \Pi_0 \arrow[u, hook] \arrow[r, hook] \& |[alias=s205gzero]| \mathcal{G}_0 \arrow[u, hook] \&
\end{tikzcd}},

where the first row is formed of profinite groups, the second of discrete groups. We will note that \(\mathcal{G}\) is identified with a quotient of the profinite completion of \(\mathcal{G}_0\), but

\scanpage{206}

I do not require for the moment that it be isomorphic to this completion itself. \unsure{Thus} \unsure{idem|item} for $\mathcal{T}$ and $\mathcal{T}_0$. The group $\mathcal{T}_0\inserted{=\mathcal{G}_0/\Pi_0}$ plays the role of a discrete \inserted{``direct''} Teichmüller group, the group $\mathcal{T}$ that of a completion of a Teichmüller group, the group $\Gamma=\mathcal{M}/\mathcal{G}$ plays the role of a profinite Galois group. \deleted{The \unsure{situation}} I am interested mainly in two extreme situations

a) $\mathcal{G}_0=\Pi_0$ whence $\mathcal{G}=\Pi$ --- situation coming from a smooth relative \unsure{$1$-multiplicity} over a field \deleted{$k$} extension of finite type of $\mathbb{Q}$, or more generally \deleted{from a scheme}\inserted{over a} scheme \inserted{(or a multiplicity)} \unsure{integral} of finite type over $\mathbb{Q}$, \deleted{\unsure{$\Gamma$|$k$} \unsure{in fact}}\inserted{$\mathcal{N}\unsure{\xrightarrow{\sim}|\simeq}\Gamma$ is} \deleted{\unsure{is}} the fundamental group of the said \deleted{\unsure{this}} scheme or of the said multiplicity,

b) $\mathcal{G}_0/\Pi_0\simeq$ Teichmüller group of $\Pi_0$ \inserted{or of finite index in this Teichmüller group} --- in this case, one is interested mainly in a ``modular'' situation, where

\scanpage{207}
\mspage{751}

$\mathcal{N}$ is interpreted as the fundamental group of \unsure{their|the} appropriate modular multiplicity, $\mathcal{T}$ as its ``geometric'' part, $\mathcal{T}_0$ as its ``transcendental'' part, and $\Gamma$ as $\Gamma_{\unsure{\overline{\mathbb{Q}}_0|\overline{\mathbb{Q}}}/\mathbb{Q}}$ ....

We make on (32) the hypothesis that the cyclotomic character on $\mathcal{M}$ (operating on $\Pi$ by automorphisms of $L$-structures) is trivial on $\mathcal{G}$, a fortiori on $\Pi$ (thus
\[
(33)\qquad \Pi=\Pi^{+},\quad \mathcal{G}=\mathcal{G}^{+}\qquad ),
\]
i.e. that it factors as
\diagram{(34)\qquad
\begin{tikzcd}[ampersand replacement=\&, column sep=large]
\Gamma \arrow[r, "\chi_{\Gamma}"] \& \widehat{\mathbb{Z}}^{*}
\end{tikzcd},
}
\uline{whose image is an open subgroup.}

If $(G_{\mathbb{Q}},\chi_G)$ is as above (page 748), we are interested in the homomorphisms (27)
\diagram{
\begin{tikzcd}[ampersand replacement=\&, column sep=large]
\mathcal{M} \arrow[r] \& G(\mathcal{A}_{\mathbb{Q}})
\end{tikzcd}
}
satisfying the conditions $1^{\circ})$ $2^{\circ})$ $\deleted{\unsure{2}}\inserted{3}^{\circ})$ $4^{\circ})$ as above, but where $3^{\circ})$ is \unsure{formulated} in the slightly stronger form

($3^{\circ}$ strong) (Compatibility with Teichmüller discretification $\mathcal{G}_0$) : We have a \unsure{factoriza-}

\scanpage{208}

tion $\mathcal{G}_0\longrightarrow G(\mathbb{Q})$, making \unsure{thus} the diagram commutative
\diagram{(35)\qquad
\begin{tikzcd}[ampersand replacement=\&, column sep=large, row sep=large]
\mathcal{G}_0 \arrow[r, hook] \arrow[d] \& \mathcal{M} \arrow[d] \\
G(\mathbb{Q}) \arrow[r, hook] \& G(\mathcal{A}_{\mathbb{Q}})
\end{tikzcd}.
}

Structures of the form (32) probably deserve a name, I am going to call \unsure{them} --- if only provisionally, \uline{Galois-Teichmüller \unsure{structures}.}
\margin{Leaving it open later to introduce more draconian conditions on \unsure{this structure}}

If $G$ is an affine algebraic group over $\mathbb{Q}$, a homomorphism $\mathcal{M}\unsure{\longrightarrow|{-}}G(\mathcal{A}_{\mathbb{Q}})$ satisfying the conditions $1^{\circ})$ $2^{\circ})$ $3^{\circ})$ \inserted{(strong)} $4^{\circ})$ is called a ``Lie homomorphism'' of the Galois-Teichmüller structure.

I will henceforth denote by $\underline{\mathcal{M}}_{\mathbb{Q}}$ \inserted{(or \unsure{simply} $\underline{\mathcal{M}}$ ?)} the group scheme \inserted{$G_{\mathbb{Q}}$} over $\mathbb{Q}$. \deleted{\unsure{defined}}

I denote by \deleted{$\unsure{\mathcal{G}|\mathcal{M}}$} $\underline{\mathcal{G}}_{\mathbb{Q}}$ the subgroup scheme of $\underline{\mathcal{M}}_{\mathbb{Q}}$ generated by the

\scanpage{209}
\mspage{752}

image of $\mathcal{G}_0\longrightarrow\deleted{\unsure{\mathcal{M}|G}}G(\mathbb{Q})$, by $\underline{\Pi}_{\mathbb{Q}}$ the one generated by the image of $\Pi\longrightarrow G(\mathbb{Q})$

We thus have \deleted{\unsure{a}} subgroup schemes
\[
(36)\qquad \underline{\Pi}_{\mathbb{Q}}\subset
\deleted{\underline{\unsure{\mathcal{M}}}}\inserted{\underline{\mathcal{G}}}_{\mathbb{Q}}
\subset\underline{\mathcal{M}}_{\mathbb{Q}}
\]
\uline{normal} in $\underline{\mathcal{M}}_{\mathbb{Q}}$, giving rise to quotient group schemes
\[
(37)\qquad
\left\{
\begin{aligned}
&\underline{\mathcal{T}}_{\mathbb{Q}}\simeq
\underline{\mathcal{G}}_{\mathbb{Q}}/\underline{\Pi}_{\mathbb{Q}},\quad
\underline{\Gamma}_{\mathbb{Q}}\simeq
\underline{\mathcal{M}}_{\mathbb{Q}}/\underline{\mathcal{G}}_{\mathbb{Q}},\\
&\qquad\underline{\mathcal{N}}_{\mathbb{Q}}\simeq
\underline{\mathcal{M}}_{\mathbb{Q}}/\underline{\Pi}_{\mathbb{Q}}
\end{aligned}
\right.
\]
--- giving rise to four exact sequences of group schemes, the last of which is
\diagram{(38)\qquad
\begin{tikzcd}[ampersand replacement=\&, column sep=large]
1 \arrow[r] \& \underline{\mathcal{T}}_{\mathbb{Q}} \arrow[r] \&
\underline{\mathcal{N}}_{\mathbb{Q}} \arrow[r] \& \underline{\Gamma}_{\mathbb{Q}} \arrow[r] \& 1
\end{tikzcd}.
}

\deleted{I \unsure{am going to|see} \unsure{reflect|recall}}\inserted{Let us note \unsure{that}} the character $\chi_G:G_{\mathbb{Q}}\longrightarrow \unsure{\mathbb{G}_m|G_m}$ is an epimorphism, and factors through $\underline{\Gamma}_{\mathbb{Q}}$ as
\diagram{(39)\qquad
\begin{tikzcd}[ampersand replacement=\&, column sep=large]
\underline{\Gamma}_{\mathbb{Q}} \arrow[r, "\chi_{\underline{\Gamma}}"] \& \unsure{\mathbb{G}_{m\mathbb{Q}}|G_{m\mathbb{Q}}}
\end{tikzcd}.
}